According to variation theory \autocite[See][]{NCOL}, learning an educational
objective requires the learner to discern all critical aspects of that
objective. Misconceptions occur when the learner cannot yet discern one or
more critical aspects; documented misconceptions therefore help us identify
what those aspects are.

In this work (in progress), we apply this approach---developed for
introductory programming in our companion paper
\autocite{VTProgMisconceptions}---to document preparation with LaTeX: the
compile model (source to output), error handling, and semantic markup.
LaTeX is an interesting case for two reasons. First, the learner arrives
with a strong prior model---the WYSIWYG word processor---against which
LaTeX's aspects must be discerned; the only controlled comparison found
shows LaTeX users slower and more error-prone than word-processor users on
plain text while enjoying their tool more \autocite{Knauff2014Latex}.
Second, the education literature on LaTeX-specific misconceptions appears
to be essentially empty, so this catalogue must be built from adjacent
literatures and our own course data. We treat documented difficulties as
phenomenographic observations, analyse them through the lens of variation
theory to identify critical aspects, and outline patterns of variation to
teach learners to discern them.
